\chapter{Conclusion}\label{ch:conclusion}

This thesis asked whether safety in a market with adaptive agents requires knowing who the agents are. For the markets studied here it does not, and the reason is structural rather than empirical. Prices reveal the agents only through a merged measure of their memories and weights, and recovering the agents from that measure costs capital of order $\delta^{-4}$, experiments of order $(\Delta\eta)^{-(4K-2)}$, and with general memory an accuracy that improves only logarithmically (Chapter~\ref{ch:p1}). Safety depends on the agents only through their total feedback weight. That one number is certified at the parametric rate up to a single logarithm under a memory bound, from passive data at a cost of order one over the stability margin, and not at all without a memory bound; a guard built on the certificate is the largest safe admission rule (Chapters~\ref{ch:p2} and~\ref{ch:p3}). The directions in which prices are least informative about the agents are directions in which the risk does not move. Certify, don't identify.

\paragraph{What the evidence supports.} The identification half holds in data: prices and order flow leave the agents' memories essentially unrestricted, on new symbols and new days, while disclosed trader-class positions restrict them. The certification half holds in a narrower form. With enough days, public data decide the sign of the stability margin in new markets whose agents they do not identify; a lower margin forecasts next-day variance beyond the standard measure of self-excitation; and in thin books the summed response forecasts the size of the next hour's adverse excursion beyond a full volatility baseline. But the decided margins rest mainly on price reversion, the markets studied were far from the boundary, the margin does not forecast adverse moves or rank damage on stress days, and a guard estimated from public data does no better than a volatility limit. Of the 23 primary hypotheses of the thesis's registered tests, 11 passed, and of the 25 of Paper~2's program, 15 passed. These are the results of one researcher's program on public data, reported with their failures, not a settled empirical case.

\paragraph{AI agents.} The question has become practical. When trading and lending decisions are delegated to AI agents built on a few shared models, the risk supervisors fear is that the agents act alike, and the monoculture proposition makes that precise: a shared model lowers the margin in proportion to its adoption share and multiplies the shocks the loop amplifies. In simulation with learning agents, the margin's trend, computed from public data alone, warned about two weeks before instability with no false alarm, while a threshold computed from a model of the agents erred by nearly a factor of two; with language-model desks it again warned before every crisis with no false alarm, but only three to four days ahead. A supervisor does not need to read the agents' code or predict how they learn; it needs a bound on their memory, renewed as often as they change, and a trigger on the loop. Test~B2 showed that the trigger's specificity survives language-model desks and its lead does not; Test~H1 will say whether real accounts on a live venue have the distinct memories the theory assumes.

\paragraph{The policy message.} Supervisors who want to control transient risk should ask for the information that certifies safety, a bound on memory and on total sensitivity, renewed as often as the agents learn, rather than for the identities or internals of the agents. Identity-based regimes serve other purposes, such as enforcement and attribution, but in the model of this thesis they are neither necessary nor efficient for safety. Passive observation of prices delivers neither kind of knowledge reliably when markets are far from the boundary; disclosure, or an outside shock that acts as an experiment, is what makes certification attainable in practice.

\section{Open problems}\label{sec:open}

\begin{enumerate}
\item \emph{Nonlinear questions.} Theorem~\ref{thm:dich}(iv) gives the exact cost of every linear question, and Proposition~\ref{prop:nonlinear} settles questions that depend on finitely many taps, such as the peak. Extend the exact theory to questions that depend nonlinearly on infinitely many taps, and make the link to the $1/\log N$ rate of identification (Corollary~\ref{cor:cases}(e)) a theorem rather than a reading of Chapter~\ref{ch:p1}'s proof.
\item \emph{Constants.} Theorem~\ref{thm:philog} gives explicit constants, but its lower constant is far below its upper one ($\PhilogLowB$ against $\PhilogUpB$ at $\bar r=0.9$). Find the limit of $\omega_L(\delta)/(\delta\sqrt{\log(1/\delta)})$; the computed brackets sit near the upper constant.
\item \emph{Minimal disclosure for safety.} Characterize the least informative disclosure under which a given safety question becomes certifiable from finite data. Theorem~\ref{thm:dich}(iii) says a memory bound is necessary for the summed response; is it sufficient in a sense that can be made operational?
\item \emph{Strategic agents.} A continuous throttle with a buffer makes every strategy profile safe and has equilibria (Proposition~\ref{prop:throttle}), and agents cannot move the certificate within the memory bound without moving the margin (Theorem~\ref{thm:nogaming}). Open: whether equilibrium sensitivity is lower with the guard than without it, and how to supervise agents that can switch their response abruptly at a time of their choosing.
\item \emph{Observational certification.} Develop finite-sample certificates from observational order-book data with dependence and non-stationarity, replacing the bootstrap diagnostics of Chapter~\ref{ch:p2}. Test~E1 found passive five-minute data far from the reset contract of Theorem~\ref{thm:dich}; natural experiments (Proposition~\ref{prop:natural}) are the nearest substitute, and their sampling theory under dependent shocks is open. Theorems~\ref{thm:passivecert}, \ref{thm:passivemulti} and~\ref{thm:uniformcert} settle the cost of passive certification for one closed-loop root and, uniformly as the margin vanishes, for any number of cohorts with finite memory; Theorem~\ref{thm:hidden} shows that without a memory bound no passive certificate exists. The pooled price-and-flow estimator that the tests use attains the same order (Theorem~\ref{thm:pooled}). What remains is a version uniform over geometric memory structures and over the pooled estimator's nuisance parameters, and the gap between the moving-average displacement the tests use and the distance from the fundamental value the theory uses.
\item \emph{The evidence.} Run Test~E2 after 2 February 2027; measure the feedback term itself, not only the margin, from disclosure or outside shocks; and extend the cohort model to contrarian agents, which the exploration finds to be common.
\end{enumerate}
